\chapter{MLOps, Observability, and Production Readiness}
\label{ch:mlops}

A trained model that cannot be retrained, versioned, or monitored is a liability rather than an asset. This chapter covers the infrastructure surrounding the models: experiment tracking, model promotion, retraining triggers, the supervisor feedback loop, and observability.

\section{MLOps Lifecycle Overview}

The MLOps layer turns the models of Chapter~\ref{ch:training} from static artifacts into a managed lifecycle with four stages: \textbf{train} (produce a candidate with logged hyperparameters and metrics), \textbf{promote} (replace production only if performance is equal or better), \textbf{serve} (the streaming pipeline loads the promoted model from a shared volume mount, requiring no image rebuild), and \textbf{monitor} (observe the production score distribution and trigger a new cycle when drift or feedback indicates degradation).

The lifecycle is largely automated. The single human touch point is the supervisor confirming or rejecting alerts in the dashboard, which feeds one of the three retraining triggers. The promotion gate in particular has no manual approval step --- a deliberate scope decision discussed below.

\section{Experiment Tracking with MLflow}

MLflow provides the tracking layer: every training job, whatever triggered it, logs to the same server, creating a single lineage of all model versions. The server runs as a Docker container with a flat-file backend (\code{mlruns/} mounted as a volume) rather than a database-backed store. For a single-user deployment running one job per week, this avoids the operational overhead of a dedicated metadata database; a production deployment serving concurrent training jobs would migrate to a PostgreSQL-backed store.

Each execution of \code{src/training/train\_all.py} produces \textbf{two independent runs} under the experiment \code{amen-anomaly-detection}, one per model. They are independent rather than nested because the two models have different architectures, hyperparameters, and metrics --- nesting would conflate unrelated configuration spaces. Each run logs parameters (architecture dimensions, learning rate, batch size, epochs, sequence length), metrics (validation AUC-ROC against the oracle labels), and artifacts (the trained \code{.pt} file, the fitted scaler, and the reference score distributions used by drift detection).

\section{Model Promotion Gate}

After training, the system must decide whether the candidate should replace the model currently serving. Each model type is evaluated independently: the script reads \code{models/production\_metrics.json}, which stores the production AUC per type, and promotes the candidate if its AUC is $\geq$ the production value --- writing its \code{.pt} file, scaler, and reference distribution to the shared model volume and updating the baseline.

The comparison uses $\geq$ rather than $>$ deliberately: a model retrained on fresher data with an identical AUC is preferable to one trained on older data, because the feature distributions it has seen are more representative of current traffic. On the very first run no baseline file exists; the script detects this and bootstraps by promoting unconditionally.

Two limitations are worth stating plainly. First, \textbf{no human approval step}: a model passing the AUC threshold is promoted immediately, whereas a production banking environment would require a compliance or risk officer to review model changes --- an Airflow task pausing the DAG until confirmation is a straightforward extension. Second, \textbf{single-metric comparison}: the gate evaluates AUC only, where a more robust gate would add F1 at the operating threshold, per-scenario recall to catch regressions on specific anomaly types, and inference latency to guard against deploying a model that is accurate but too slow for real-time scoring.

\section{Orchestration with Apache Airflow}

Airflow runs with a \code{LocalExecutor}, executing tasks in the scheduler process rather than distributing them --- sufficient for two DAGs at low concurrency. Migrating to \code{CeleryExecutor} is a configuration change requiring no DAG modifications, only worker containers and a message broker. The Airflow metadata tables share the application PostgreSQL instance to conserve memory, and a dedicated \code{Dockerfile.airflow} installs only orchestration packages --- excluding PyTorch, FastAPI, and SHAP --- keeping the image roughly 2\,GB smaller than the training image.

Beyond the nightly profile DAG of Chapter~\ref{ch:streaming}, the \code{weekly\_retrain} DAG contains three tasks: \emph{validate data freshness} (short-circuit if the \code{transactions} table has received no new events since the last run, avoiding wasted compute on an identical training set), \emph{retrain models} via a \code{DockerOperator}, and \emph{verify promotion} (confirm \code{production\_metrics.json} was updated).

The \code{DockerOperator} was chosen over a Bash or Python operator for two reasons. \emph{Dependency isolation}: training requires PyTorch, which is large and unnecessary in the scheduler. \emph{Reproducibility}: the training image is pinned as \code{amen-anomaly-training:latest} with \code{force\_pull=False}, so the same image serves scheduled, drift-triggered, and feedback-triggered retraining, and is rebuilt explicitly when training code changes rather than pulled on every run. The model volume is mounted into the training container at the same path used by the Scoring Service, so a promoted model is immediately available on the next restart --- no copy or deployment step.

\section{Retraining Triggers}

Three independent paths can initiate retraining, and all converge on the same \code{weekly\_retrain} DAG so that every job follows the same training, promotion, and artifact-logging pipeline.

\textbf{Scheduled.} A weekly cron in the DAG definition provides a baseline cadence, refreshing the models with recent data even when no anomaly is detected.

\textbf{Drift-triggered.} During training, the script saves the validation-set score distributions as reference arrays logged to MLflow. In production, the Scoring Service exposes the fraction of events scoring above $T_1$ as a Prometheus gauge, and an alert rule fires when that fraction exceeds 15\% for a sustained hour:

\begin{center}
\code{alert: ScoreDistributionDrift \quad expr: score\_above\_t1\_ratio > 0.15 \quad for: 1h}
\end{center}

The threshold and window were chosen empirically: under normal operation fewer than 5\% of events exceed $T_1$, so a sustained elevation above 15\% indicates a material distribution shift rather than a single burst of legitimate anomalies. Figure~\ref{fig:drift} shows the resulting chain.

\begin{figure}[htbp]
\centering
\begin{tikzpicture}[node distance=0.7cm and 0.9cm, every node/.style={font=\small}]
  \tikzstyle{box}=[draw, rounded corners=2pt, align=center, minimum height=0.9cm, inner xsep=5pt]
  \node[box] (p) {Prometheus\\{\tiny alert rule}};
  \node[box, right=of p] (am) {AlertManager\\{\tiny port 9093, 24h dedup}};
  \node[box, right=of am] (wh) {Webhook relay\\{\tiny FastAPI, auth}};
  \node[box, right=of wh] (af) {Airflow REST\\{\tiny basic auth}};
  \node[box, below=of am, xshift=1.2cm] (dag) {\code{weekly\_retrain} DAG};
  \draw[-Stealth] (p) -- (am);
  \draw[-Stealth] (am) -- (wh);
  \draw[-Stealth] (wh) -- (af);
  \draw[-Stealth] (af.south) |- (dag.east);
\end{tikzpicture}
\caption{Drift-triggered retraining chain. The webhook relay exists because AlertManager cannot natively authenticate against Airflow's REST API.}
\label{fig:drift}
\end{figure}

AlertManager deduplicates with a 24-hour \code{repeat\_interval} to prevent repeated triggers from one drift episode, then routes to a webhook receiver. That receiver --- a single-endpoint FastAPI service on a minimal \code{python:3.11-slim} image ($\sim$150\,MB) --- exists purely because AlertManager's native webhook format does not support the HTTP Basic Authentication that Airflow's REST API requires. It performs the authentication and forwards the trigger with \code{conf=\{"trigger": "drift\_alert"\}}.

\textbf{Feedback-triggered.} Described in the next section.

\section{Supervisor Feedback Loop}

The feedback loop is the single human touch point in the lifecycle. When the pipeline produces an ALERT or BLOCK, the branch supervisor reviews it in the dashboard and either confirms or rejects it. The verdict is submitted to \code{POST /api/v1/feedback} on the FastAPI service, which writes it to the \code{evaluation\_labels} table with a safe foreign key to the original transaction.

These labels serve two purposes. They provide ground truth for evaluating model performance on real decisions --- a capability that becomes meaningful once the system runs on real transactions rather than synthetic events with oracle labels. And they feed a daily \textbf{feedback divergence DAG}, which queries what fraction of ALERT and BLOCK events were rejected over a rolling seven-day window and triggers \code{weekly\_retrain} if the rejection rate exceeds 50\% with at least 10 reviewed events.

The 50\% threshold reflects a simple principle: if more than half of the system's high-confidence alerts are being overridden by human reviewers, the models are producing more noise than signal at the upper tiers. The 10-review minimum prevents a single rejected alert early in a deployment from triggering an unnecessary cycle.

\begin{table}[htbp]
\centering
\small
\caption{Retraining trigger comparison}
\begin{tabular}{@{}llll@{}}
\toprule
\textbf{Trigger} & \textbf{Source} & \textbf{Condition} & \textbf{Frequency} \\
\midrule
Scheduled & Airflow cron & Weekly & Fixed \\
Drift & Prometheus alert & $>$15\% above $T_1$ for 1h & On demand \\
Feedback & Supervisor reviews & $>$50\% rejection, $\geq$10 reviews & Daily check \\
\bottomrule
\end{tabular}
\end{table}

The \code{conf} parameter passed at trigger time records which path initiated a run, giving traceability in the Airflow logs without separate DAG definitions per trigger type.

\section{Observability and Monitoring}

Each streaming service embeds a \code{ServiceMetrics} class exposing two Prometheus instruments: a \textbf{Counter} of events processed, labeled by service and outcome, and a \textbf{Histogram} of per-event processing latency enabling percentile queries in Grafana. Services expose metrics on dedicated ports (enrichment 9100, scoring 9101, decision 9102), and a Kafka exporter (9308) provides broker-level metrics including \emph{consumer group lag} --- the gap between the latest produced and latest committed offset --- which is the primary indicator of pipeline backpressure.

\textbf{Failure domain separation.} A monitoring system must not share a failure domain with the service it monitors. If the enrichment service and its health reporter both depended on Redis, a Redis outage would simultaneously break the service and suppress the alert reporting the breakage. This ruled out Redis-based heartbeats early in the design. The chosen architecture has Prometheus pull metrics over HTTP from each service's embedded endpoint; the only shared dependency is the network layer itself, an acceptable residual risk since no monitoring architecture can report a network outage from inside the affected network.

\textbf{Structured logging.} A \code{JsonFormatter} class emits one JSON object per log line --- timestamp, service, level, message, plus contextual extras such as \code{event\_id}, \code{client\_id}, \code{error\_type}, and \code{alert\_tier}. A \code{setup\_logger(service\_name)} factory configures the standard logging module with this formatter, replacing all \code{print()} statements from earlier versions. Structured logs enable filtering by service, grouping by error type, or correlating events across services by \code{event\_id} without fragile regular-expression parsing of free text.
